\section{Tree Operations}

\subsection{Implement \emph{depths}}
Below my implementation of the depths can be seen, which seems to work.
\begin{lstlisting}[language={futhark}]
def depths (steps: []step) : [](i64, i32) =
  let (arr, values, keep) = unzip3 <| map (\s -> 
    match s
    case #u -> (-1, -1, false)
    case #d x -> (1, x, true)
  ) steps
  let inc_scan_arr = scan (+) 0 arr
  let exc_scan_arr = map2 (\a b -> a - b) inc_scan_arr arr
  let (depth, value, _) = zip3 exc_scan_arr values keep |> filter (\(_, _, b) -> b) |> unzip3
  in zip depth value
\end{lstlisting}


Here we have to convert the steps into items of 1 for down and -1 for up
operations. Then we can do an exclusive scan and at last remove all the
elements which originally corresponded to up items as these do not add new
elements to the graph.

If we let N denote the length of steps, then we have a map, a scan, another map
and at last a filter, which all work on an array of the full length N. The
filter and the scan operation is what is the limiting factor giving us work of
O($N$) and span of O($\lg N$).

Any sequential computation would have to touch all elements at least once, and
as such this implementation is work efficient.

\subsection{Implement \emph{parents}}
Below is my implementation of the parents function.
\begin{lstlisting}[language={futhark}]
def parents (D: []i64) : []i64 =
  let n = length D
  let depth_tuple = zip (D :> [n]i64) (iota n)
  let parents = loop parents = [0] for i' < n - 1 do 
    let i = i' + 1
    let search_d = D[i] - 1
    let search_ds = map (\i -> depth_tuple[i]) (iota i)
    let (_, p) = reduce_comm (\(d1, i1) (d2, i2) -> 
      if (d1 == d2) && (d1 == search_d) then 
        if i1 > i2 then (d1, i1)
        else (d2, i2)
      else if d1 == search_d then (d1, i1)
      else if d2 == search_d then (d2, i2)
      else (-1,-1)
    ) (0, -1) search_ds
    in parents ++ [p]
  in parents
\end{lstlisting}
We initialize the parents with the root being itself, and then otherwise we
search to the left of us for items with a value less than the current index,
where we choose the higher index in case of a tie. 

The algorithm runs a sequential loop over all elements (except the root).
Inside this loop it runs a map and a reduce\_comm. The reduce\_comm only runs
on the left items, meaning that the length that is searched is a reduced subset
of the full list. However, as an upper bound this would still be over the full
array. As such inside the loop we do O($N$) work with a span of O($\lg N$). We
do this within a loop giving total work of O($N^2$) and span O($N \lg N$).

We are able to compute the parent vector by a single traversal of the original
tree and as such I would not state that this algorithm is work efficient. 

\subsection{Implement \emph{subtree\_sizes}}
Below my implementation of the \emph{subtree\_sizes} can be seen.
\begin{lstlisting}[language={futhark}]
def subtree_sizes [n] (steps: [n]step) : []i64 =
  let (D, V) = depths steps |> unzip 
  let P = parents D
  let n = length D
  let max_depth = reduce i64.max 0 D
  let (res, _) = loop (res, cur_depth) = (copy V, max_depth) while cur_depth > 0 do
      let indices = filter (\i -> D[i] == cur_depth) (iota n)
      let values = map (\i -> res[i]) indices
      let parent_indices = map (\i -> P[i]) indices
      in (reduce_by_index res (+) (0) parent_indices values, cur_depth - 1)
  in map i64.i32 res
\end{lstlisting}

This approach uses the naive bottom up iteration. Here we have a call to the
reduce function, followed by a sequential loop, which runs over the depth of
the tree. This loop calls a filter, map and reduce\_by\_index. 


